\documentclass[11pt]{article}

\usepackage[margin=1in]{geometry}
\usepackage[utf8]{inputenc}
\usepackage[T1]{fontenc}
\usepackage{lmodern}
\usepackage{microtype}
\usepackage{graphicx}
\usepackage{booktabs}
\usepackage{amsmath, amssymb}
\usepackage{array}
\usepackage{multirow}
\usepackage{float}
\usepackage[hidelinks]{hyperref}
\usepackage[numbers,sort&compress]{natbib}
\usepackage{xcolor}
\usepackage{caption}
\usepackage{lineno}
\linenumbers
\captionsetup{font=small,labelfont=bf}

\title{\textbf{Biological Continued Pretraining Reshapes the Capability\\
Profile of a Foundation Model Without Catastrophic Forgetting}}

\author{Liang Wang\thanks{Corresponding author. Email: wangliang.f@gmail.com}\\
\small School of Artificial Intelligence and Automation\\
\small Huazhong University of Science and Technology\\
\small Wuhan, Hubei 430070, P.R. China}
\date{July 2026}

\begin{document}
\maketitle

\noindent\textbf{Highlights}
\begin{itemize}\itemsep2pt
\item A training-free re-analysis of one 26B MoE lineage isolates the effect of biological
continued pretraining (CPT) from tokenizer changes and from fine-tuning.
\item Biological CPT does not cause catastrophic forgetting; it \emph{raises} general knowledge
(MMLU $+13$ pts) and code generation (MBPP pass@1 $0.33\!\to\!0.63$).
\item CPT also makes chain-of-thought reasoning 41\% shorter and near-backtrack-free while
preserving accuracy --- an effect invisible to accuracy metrics.
\item A consistent CPT-lifts / SFT-narrows division of labor recurs across four axes, reframing
biological sequence as structured scientific data that reshapes a model's capability profile.
\item A controlled seven-model experiment (varying only the CPT data mixture) confirms this
causally: raising biological share to 50\% preserves general ability (MMLU $-0.6$ pt) while
biology rises monotonically --- data mixture is a general-preserving lever on capability.
\item A small-model, matched-compute causal probe (GPT-2, 3 seeds) finds that biological CPT
\emph{reproducibly lifts} paraphrase detection (PAWS-en, $+3.1$ pt) --- the direct structural
counterpart of protein-homology transfer --- alongside a real cost to long-range discourse
modeling, an honestly reported partial, not uniform, transfer.
\end{itemize}

\bigskip
\noindent\fbox{\parbox{\linewidth}{%
\textbf{The Bigger Picture.}\\[2pt]
Adapting a general-purpose AI model to a specialized domain --- here, the language of DNA and
proteins --- is usually assumed to come at a cost: teach it biology and it forgets how to reason
about everything else. This ``no free lunch'' intuition shapes how practitioners budget compute
and whether they attempt domain adaptation at all. We test the assumption directly, and without
running any new training, by comparing three snapshots of the same model taken before and after
biological training. The result overturns the intuition: feeding the model raw biological
sequence made it \emph{better} at general knowledge, at writing code, and even changed \emph{how}
it reasons --- producing shorter, more decisive chains of thought without losing accuracy. The
gains appear during the sequence-pretraining stage and are partly given back during
task-specific fine-tuning, revealing that the two stages play complementary rather than
interchangeable roles. This suggests a broader principle for data-centric AI: structured
scientific data --- biological sequence today, and by extension code, mathematics, and
chemistry --- is not merely knowledge to be absorbed but a lever that reshapes what a foundation
model can do. A controlled follow-up, in which we retrain the model seven times while changing
only the data mixture, confirms the lever directly: more biological data buys more scientific
capability at almost no cost to general ability. A final, smaller-scale experiment asks the
question in reverse: does biological training feed anything back into ordinary language
understanding? Under a causal design that holds compute exactly fixed, it does --- reproducibly
improving a model's ability to detect when two sentences paraphrase each other --- though not
without cost elsewhere, a nuance we report honestly rather than smooth over.
}}
\bigskip

\begin{abstract}
\noindent It is widely assumed that continued pretraining (CPT) on a narrow, out-of-distribution
corpus such as raw biological sequence must trade away a general-purpose model's broad
competence --- the ``alignment tax'' or catastrophic-forgetting intuition. We test this
directly, without any new training, by re-analyzing three checkpoints from a single lineage
of a 26B-parameter Mixture-of-Experts model (Gemma-4-26B-A4B): the instruction-tuned base,
the same model after biological CPT (8.7B tokens of DNA, protein, and biomedical text), and
after subsequent supervised fine-tuning (SFT). Across three independent capability axes ---
general knowledge/reasoning (MMLU, ARC, HellaSwag), code generation (MBPP), and biomedical
knowledge (BixBench) --- we find that biological CPT does \emph{not} degrade the model; it
\textbf{lifts} it: MMLU $+13$ points, MBPP pass@1 nearly doubles ($0.33\!\to\!0.63$), and
BixBench discrimination rises sharply (MCC $0.23\!\to\!0.92$). The single measured regression
is truthfulness (TruthfulQA $-8.8$ points), a small and interpretable domain drift. A clean
vocabulary-expansion ablation ($<\!0.4$ pt on every general metric) confirms the gains are
attributable to CPT, not tokenizer changes. Crucially, subsequent SFT \emph{narrows} the model
back: all three axes fall to near-base levels, revealing a consistent division of labor ---
\textbf{CPT re-organizes and lifts the shared capability substrate; SFT cashes it out onto
target tasks}. We argue this reframes biological sequence not as a competitor for a foundation
model's capacity but as a form of structured scientific data that reshapes its capability
profile, and that CPT and SFT should be budgeted as complementary rather than substitutable
stages. We then confirm the mechanism causally: continuing pretraining of the same base seven
times under one fixed recipe, varying only the data mixture ($0/5/20/50\%$ biological share plus
a protein/DNA composition ablation), shows that raising the biological share to $50\%$ leaves the
general axis essentially unchanged (MMLU $-0.6$ pt, GSM8K $-1.2$ pt) while biological capability
rises monotonically (remote-homology MCC flips negative to positive; BixBench saturates), and
the \emph{type} of scientific data selects which capability is amplified. Mixture ratio is thus
a controllable, general-preserving lever on the capability profile. Finally, we ask the reverse
question at a scale with capability headroom the 26B model lacks: does biological CPT feed
anything back into ordinary language understanding? A matched-compute continuation experiment on
GPT-2 (124M), holding token budget, steps, and optimizer identical between a biology-continued
and a text-continued arm across three seeds, finds a real, reproducible transfer to the direct
structural counterpart of protein-homology detection --- paraphrase detection (PAWS-en,
$+3.1$ points, consistent in sign across all three seeds) --- alongside a genuine cost to
long-range discourse modeling (LAMBADA, $-14.6$ points), a partial rather than uniform
dissociation that we report honestly. All checkpoints, evaluation code, and per-example outputs
are public.
\end{abstract}

\medskip
\noindent\textbf{Keywords:} foundation models, continued pretraining, catastrophic forgetting,
data-centric AI, mixture-of-experts, capability transfer, structured scientific data.

% ===================================================================
\section{Introduction}
\label{sec:intro}


The dominant intuition about continued pretraining (CPT) on a narrow domain corpus is
subtractive: pouring out-of-distribution tokens --- raw protein and DNA sequence, which share
almost no surface statistics with natural language --- into a general-purpose model is expected
to erode its broad competence, a phenomenon variously described as the \emph{alignment tax} or
\emph{catastrophic forgetting} \citep{luo2023forgetting, ibrahim2024cpt}. Under this view,
domain adaptation is a trade: one buys biological ability at the cost of general ability, and
the practitioner's job is to minimize the damage.

This paper argues that, for a modern Mixture-of-Experts foundation model, the intuition is
wrong --- or at least incomplete. We show, by \emph{re-analysis of existing checkpoints and
without any new training}, that biological CPT on a 26B-parameter model does not merely avoid
catastrophic forgetting; it \emph{improves} the model on capability axes that have nothing to do
with biology, including general knowledge (MMLU~\citep{hendrycks2020measuring}) and code
generation (MBPP). We then show that
the subsequent supervised fine-tuning (SFT) stage, far from adding capability, \emph{narrows}
the model back toward its base competence on held-out axes. The picture that emerges is a clean
division of labor between the two training stages, and it suggests a reframing of what
biological data \emph{is} to a foundation model.

Our substrate is a Mixture-of-Experts model~\citep{shazeer2017sparse, fedus2022switch,
jiang2024mixtral} (Gemma-4-26B-A4B~\citep{gemmateam2025}); the analysis is deliberately
training-free, re-using an existing CPT/SFT lineage~\citep{ibrahim2024cpt} so that every
capability difference is attributable to a single stage. This complements router-level
interpretability of the same model family~\citep{wang2026omnigene, park2024moeinterpret,
kanai2025moeinterp} by asking not \emph{where} computation is organized but \emph{how} the
training data reshapes what the model can do.

\paragraph{Related work.}
Our claim sits at the intersection of three literatures. \emph{Continued pretraining and
forgetting}: domain-adaptive CPT is standard practice~\citep{ibrahim2024cpt, jindal2024balance,
jiang2024instruction}, but is widely reported to risk catastrophic forgetting of general
ability~\citep{luo2023forgetting, huang2024sssr}, motivating replay and balancing schemes. Most
of this evidence is on dense models and measures forgetting as a cost to be minimized; we instead
find a net \emph{gain} on a sparse MoE model and attribute it to the structural regularity of the
pretraining corpus rather than its factual content. \emph{Biological foundation models}: protein
and genomic models --- ESM-2~\citep{lin2023evolutionary}, ProtT5~\citep{elnaggar2021prottrans},
DNABERT-2~\citep{zhou2023dnabert2}, Nucleotide Transformer~\citep{dallatorre2023nucleotide}, the
Evo family~\citep{nguyen2024evo, brixi2026evo2}, and cross-modal LucaOne~\citep{he2024lucaone}, as
well as LLMs adapted to biology~\citep{luo2022biogpt, wang2026biopaws, wang2024llamagene} --- are
typically evaluated only on in-domain tasks, leaving open whether biological training helps or
harms general competence. \emph{Data-centric AI and scaling}: work on pretraining-data
composition and its effect on capability~\citep{gadre2021pile, hoffmann2022chinchilla} argues that
\emph{what} a model is trained on shapes \emph{what} it can do; we provide a controlled,
single-lineage instance of this principle where the manipulated variable is a scientific-sequence
corpus and the readout spans knowledge, code, and reasoning behavior. To our knowledge no prior
work reports that biological CPT \emph{improves} general and coding benchmarks, nor that it
measurably alters reasoning style.

\paragraph{Contributions.}
(i) A training-free re-analysis protocol over a single checkpoint lineage that isolates the
effect of biological CPT from tokenizer changes and from SFT. (ii) Evidence that biological CPT
lifts, rather than degrades, general and code capability, with a single interpretable regression
(truthfulness). (iii) A consistent CPT-lifts / SFT-narrows division of labor observed across four
independent evaluation axes. (iv) A data-centric reframing: biological sequence as structured
scientific data that reshapes a foundation model's capability profile. (v) A controlled
seven-model mixture experiment that turns the observational claim causal: varying only the CPT
data mixture, general ability is preserved across a $0\!-\!50\%$ biological share while biological
ability rises monotonically, and data \emph{composition} selects which capability is amplified.
(vi) A small-model, matched-compute causal probe of the reverse direction, showing that
biological CPT reproducibly improves natural-language paraphrase detection --- the direct
counterpart of the protein-homology transfer that motivates this line of work --- while honestly
reporting a real cost to long-range discourse and a scale at which character-level structural
probes are simply unmeasurable.

% ===================================================================
\section{Setup: one lineage, three checkpoints}
\label{sec:setup}

Our analysis rests on a single training lineage, so that every capability difference we
report can be attributed to a specific training stage rather than to confounds of architecture,
tokenizer, or data version.

\paragraph{Base model.} All checkpoints derive from Gemma-4-26B-A4B, a Mixture-of-Experts
transformer (128 experts per layer, top-8 routing, ${\approx}3.8$B active parameters). We use
the instruction-tuned release as the origin of the lineage; the biological CPT and SFT
checkpoints are those of OmniGene-4~\citep{wang2026omnigene}.

\paragraph{The three (four) checkpoints.} The lineage is
\[
\text{it} \;\xrightarrow{\text{vocab-expand}}\; \text{it-bio} \;\xrightarrow{\text{bio-CPT}}\;
\text{BioCPT} \;\xrightarrow{\text{SFT}}\; \text{BioCPT+SFT}.
\]
\textbf{(1) Base (it):} the original instruction-tuned model, vocabulary $262{,}144$.
\textbf{(2) Base (it-bio):} identical weights with the vocabulary expanded to $290{,}172$ by
adding $28{,}028$ biological BPE tokens (DNA/protein), but \emph{no further training}; this
serves as a clean control that isolates the cost of tokenizer expansion alone.
\textbf{(3) BioCPT:} it-bio after continued pretraining on $8.7$B tokens of mixed biological
corpus (DNA, protein sequence, and biomedical literature text), with LoRA ($r{=}64$) over
attention, MLP, and the MoE router, embeddings unfrozen; merged to full weights.
\textbf{(4) BioCPT+SFT:} BioCPT after supervised fine-tuning on biological instruction data
(a dual-head structural variant, ``v5''); merged to full weights.
We verified the lineage by direct weight comparison: it and it-bio share router/expert weights
bit-for-bit (mean abs.\ diff.\ $0$), differing only in embeddings, while BioCPT's router and
expert weights have measurably moved.

\paragraph{Evaluation.} We use the \texttt{lm-evaluation-harness} (v0.4.12) for all general and
coding benchmarks, at community-standard few-shot settings (MMLU 5-shot, ARC-Challenge 25-shot,
HellaSwag 10-shot, TruthfulQA-mc2 0-shot; HumanEval 0-shot, MBPP 3-shot). Multiple-choice tasks
are scored by loglikelihood, which is robust to a model's chat-vs-completion output format ---
a property we exploit deliberately, since our lineage mixes an instruction-tuned base with a
completion-style CPT checkpoint. Biomedical discrimination (BixBench) is likewise scored by
loglikelihood of the answer choices, and reasoning-behavior diagnostics (\S\ref{sec:results})
use a shared 8-shot chain-of-thought scaffold. All runs are on a single 96\,GB GPU in bfloat16;
the full model ($\approx$52\,GB) fits on one card, and the entire study required no training.

\paragraph{Honest scope.} Two evaluation regimes are format-sensitive and we flag them up front.
HumanEval 0-shot (bare function-completion) is unfair to instruction/SFT checkpoints and we
exclude it from conclusions, reporting MBPP 3-shot as the fair coding metric. Zero-shot protein
\emph{homology} is not measurable by any simple probe we tried (forced Yes/No loglikelihood is
near-random for all checkpoints; free generation is format-confounded); we therefore treat
sequence-homology ability separately and do not force it into the transfer analysis.


% ===================================================================
\section{Part I --- Observational re-analysis of one lineage}
\label{sec:results}

The central finding is a single pattern that recurs across every capability axis we measured:
biological CPT \emph{lifts} the model, and subsequent SFT \emph{narrows} it back
(Figure~\ref{fig:ucurve}, Figure~\ref{fig:radar}). We present the three completed axes, then the
reasoning-behavior diagnostics.

\begin{figure}[h]\centering
\includegraphics[width=0.82\linewidth]{figs/fig1_ucurve.pdf}
\caption{Every held-out capability axis follows the same trajectory: flat across the
vocabulary-expansion control (it $\to$ it-bio), a sharp rise at biological CPT (shaded), and a
fall back toward base after SFT. Raw scores; see Tables~\ref{tab:general}--\ref{tab:coding}.}
\label{fig:ucurve}
\end{figure}

\subsection{Vocabulary expansion is free}
Before attributing anything to CPT we rule out the tokenizer. The original model (it) and the
vocabulary-expanded but untrained model (it-bio) are indistinguishable on every general
benchmark (MMLU $0.646$ vs.\ $0.642$; ARC-C $0.367$ vs.\ $0.367$; HellaSwag-norm $0.488$ vs.\
$0.486$; TruthfulQA $0.533$ vs.\ $0.532$; all $\le0.4$ pt). Adding $28$K biological tokens to
the vocabulary therefore costs nothing in general competence, and any change at the BioCPT stage
is attributable to the continued pretraining itself.

\subsection{General knowledge and reasoning rise, not fall}
Table~\ref{tab:general} shows the (1) General axis. Against the intuition of catastrophic
forgetting, biological CPT \emph{improves} the model: MMLU $+13.0$ points ($0.646\!\to\!0.776$),
ARC-Challenge accuracy nearly doubles ($0.367\!\to\!0.700$), and HellaSwag-norm rises
$+36.7$ points. The one regression is TruthfulQA ($-8.8$ points), an interpretable drift on
truthfulness/anti-misinformation after heavy exposure to sequence and literature text. We note
honestly that part of the ARC/HellaSwag jump is mechanistic --- an instruction-tuned model is
weak at raw-loglikelihood multiple choice, and CPT on raw text restores base-model-like MC
behavior --- but MMLU ($+13$ pt, knowledge-heavy) is a hard gain unaffected by this. Subsequent
SFT returns all four metrics to near-base levels.

\begin{table}[h]\centering\small
\caption{(1) General. Biological CPT lifts general knowledge/reasoning; SFT narrows it back.
Loglikelihood multiple-choice; community-standard few-shot.}
\label{tab:general}
\begin{tabular}{lccccc}
\toprule
Checkpoint & MMLU (5s) & ARC-C acc (25s) & ARC-C norm & HellaSwag norm (10s) & TruthfulQA-mc2 \\
\midrule
Base (it)        & 0.646 & 0.367 & 0.403 & 0.488 & 0.533 \\
Base (it-bio)    & 0.642 & 0.367 & 0.402 & 0.486 & 0.532 \\
\textbf{BioCPT}  & \textbf{0.776} & \textbf{0.700} & \textbf{0.730} & \textbf{0.855} & 0.445 \\
BioCPT+SFT       & 0.635 & 0.382 & 0.396 & 0.486 & 0.562 \\
\bottomrule
\end{tabular}
\end{table}

\subsection{Code capability nearly doubles}
Code is a structured formal language, and the (3) Coding axis (Table~\ref{tab:coding}) shows the
same shape. On MBPP (3-shot, the format-fair metric), pass@1 rises from $0.33$ at base to
$\mathbf{0.63}$ after BioCPT, then falls to $0.35$ after SFT. (HumanEval 0-shot is excluded:
bare completion is unfair to instruction/SFT checkpoints, which score ${\approx}0$; we report
it only for completeness.) That biological-sequence CPT nearly doubles Python coding ability is
consistent with the hypothesis that structured-sequence pretraining strengthens a shared
substrate for formal/structural languages.

\begin{table}[h]\centering\small
\caption{(3) Coding. MBPP 3-shot is the fair metric (HumanEval 0-shot excluded, see text).}
\label{tab:coding}
\begin{tabular}{lcc}
\toprule
Checkpoint & HumanEval 0s (excl.) & \textbf{MBPP 3s pass@1} \\
\midrule
Base (it)        & 0.012 & 0.332 \\
Base (it-bio)    & 0.000 & 0.350 \\
\textbf{BioCPT}  & 0.439 & \textbf{0.630} \\
BioCPT+SFT       & 0.000 & 0.348 \\
\bottomrule
\end{tabular}
\end{table}

\subsection{Biomedical knowledge: the same curve}
On BixBench (biomedical true/false, loglikelihood-scored MCC), the (2) Biology axis reproduces the
pattern exactly: MCC $0.232$ at base $\to$ $\mathbf{0.924}$ after BioCPT $\to$ $0.361$ after SFT.
Protein sequence-homology, by contrast, is not measurable zero-shot by any simple probe (forced
Yes/No loglikelihood is near-random for all checkpoints; see \S\ref{sec:setup}); we read this as
evidence that sequence-homology ability is \emph{installed into representations} by CPT but only
\emph{cashed out} into task performance by SFT, consistent with prior dual-mode results on this
model family~\citep{wang2026omnigene}.

\subsection{Reasoning behavior: CPT shifts the observable reasoning trajectory}
The three axes above concern \emph{what} the model knows. We finally ask whether CPT changes the
\emph{observable trajectory} of its chain-of-thought generations, independent of accuracy --- a
behavioral phenotype of the output text, not a claim about internal reasoning process. On $200$
GSM8K problems, each sampled $k{=}5$ times under a shared 8-shot chain-of-thought scaffold, we
measure generated chain length, self-correction markers (``wait/actually/however/reconsider''),
hedging markers, and self-consistency (agreement of the final answer across the $5$ samples).
Table~\ref{tab:cot} and Figure~\ref{fig:reasoning} show a striking shift in this observable
phenotype at the BioCPT stage: the generated chain \emph{shortens by 41\%} ($108.8\!\to\!64.4$
tokens) and self-correction markers \emph{nearly vanish} ($0.370\!\to\!0.006$ backtracks per
generation), yet accuracy does not fall (it rises slightly, $0.667\!\to\!0.689$) and
self-consistency is unchanged ($0.745$). Biological CPT --- pretraining on deterministic,
low-redundancy sequence continuation --- appears to shift the model's generated text toward a
\emph{shorter, less self-revising} trajectory. Subsequent SFT reverses this: generated chains
grow past base length ($125.8$ tokens), backtracking and hedging markers return. This is the
same CPT-shapes / SFT-re-shapes signature seen on the knowledge axes, now visible in the
generated reasoning \emph{trajectory} rather than outcome --- a phenotype invisible to accuracy
alone. (it and it-bio again coincide, confirming the clean vocabulary ablation.)

\begin{figure}[h]\centering
\includegraphics[width=0.92\linewidth]{figs/fig3_reasoning.pdf}
\caption{Observable reasoning-trajectory phenotype on GSM8K. \textbf{Left:} biological CPT
shortens the generated chain-of-thought by 41\%. \textbf{Right:} self-correction markers
(backtracks) nearly vanish at BioCPT while accuracy is preserved (even slightly improved); SFT
reverses both. CPT shifts the generated trajectory toward shorter, less revision-heavy text.}
\label{fig:reasoning}
\end{figure}

\begin{table}[h]\centering\small
\caption{(4) Reasoning-behavior diagnostics on GSM8K (200 questions, $k{=}5$ samples, shared
8-shot CoT). CPT shortens the generated chain and removes self-revision markers without hurting
accuracy; SFT reverses both.}
\label{tab:cot}
\begin{tabular}{lccccc}
\toprule
Checkpoint & Chain len (tok) & Backtrack/gen & Hedge/gen & Self-consist. & Acc.\ (2nd) \\
\midrule
Base (it)        & 108.8 & 0.370 & 0.001 & 0.749 & 0.667 \\
Base (it-bio)    & 107.5 & 0.324 & 0.000 & 0.728 & 0.652 \\
\textbf{BioCPT}  & \textbf{64.4} & \textbf{0.006} & 0.000 & 0.745 & 0.689 \\
BioCPT+SFT       & 125.8 & 0.396 & 0.018 & 0.757 & 0.703 \\
\bottomrule
\end{tabular}
\end{table}

\subsection{Summary}
Across every completed axis the ordering is identical: base $<$ \textbf{BioCPT} $>$ SFT, with
BioCPT the peak (Figure~\ref{fig:radar}). Biological CPT is not a tax; it is an investment in a
shared capability substrate. SFT then trades breadth for task focus.

\begin{figure}[h]\centering
\includegraphics[width=0.60\linewidth]{figs/fig2_radar.pdf}
\caption{Capability profile of the three trained states. BioCPT (red) dominates the base model
on every axis; SFT (blue) retreats toward base on all held-out axes. The biological CPT stage,
not SFT, is what expands the profile.}
\label{fig:radar}
\end{figure}


% ===================================================================
\section{Part II --- Controlled mixture experiment}
\label{sec:phase2}

The re-analysis above is observational: it characterizes one existing lineage and cannot,
on its own, establish that the \emph{amount} and \emph{composition} of biological data are what
drive the capability shift, rather than some idiosyncrasy of that particular CPT run. We
therefore ran a controlled experiment. Starting from the same vocabulary-expanded base (it-bio),
we trained seven fresh CPT models under a single fixed recipe, varying only the training-data
mixture, and evaluated all seven on the same axes. This turns the observational claim of
\S\ref{sec:results} into a causal one: mixture ratio is a controllable lever that shapes the
capability profile.

\paragraph{Controlled CPT protocol.} Each of the seven models continues pretraining from it-bio
for exactly $100$M tokens (QLoRA, $r{=}64$ over attention, MLP, and the MoE router, with
embeddings unfrozen --- identical to the BioCPT recipe of \S\ref{sec:setup}), on a single
96\,GB GPU. The \emph{only} independent variable is the data mixture. All sources are tokenized
once into per-source pools and sampled to the target ratio, so the token budget is held at
$100$M ($\pm0.003\%$) across every run. Two groups vary orthogonally:
\emph{(i)~Mixture ratio} --- the biological share of the corpus is set to $0/5/20/50\%$ (the
remainder English web text), with the biological portion held at a fixed internal composition
($70\%$ protein, $20\%$ DNA, $10\%$ biomedical literature); \emph{(ii)~Composition} --- the
biological share is fixed at $20\%$ and its \emph{content} is varied (protein-only, DNA-only,
protein+DNA), isolating which type of scientific data drives which capability. The 20\%-mixed
model (\texttt{M2-bio20}) serves as the shared protein+DNA+literature reference point for both
groups. Evaluation reuses the general axis (MMLU 5-shot, GSM8K 5-shot) and the biology axis
(protein homology, standard and remote; BixBench), scored by loglikelihood exactly as before.

\paragraph{Result 1: general capability is preserved across the whole mixture range.}
Table~\ref{tab:p2mix} and Figure~\ref{fig:p2mix} show that raising the biological share from
$0$ to $50\%$ costs essentially nothing on the general axis: MMLU moves $0.802\!\to\!0.796$
($-0.6$ pt) and GSM8K $0.873\!\to\!0.861$ ($-1.2$ pt) --- both within run-to-run noise. There is
no mixture threshold at which the general axis collapses; even a corpus that is half
out-of-distribution biological sequence leaves general knowledge and mathematical reasoning
intact. This is the causal counterpart of the ``no catastrophic forgetting'' observation of
\S\ref{sec:results}.

\paragraph{Result 2: biological capability rises monotonically with biological share.}
Over the same sweep the biology axis climbs steadily: protein-homology discrimination (standard)
rises from MCC $0.228$ to $0.375$ ($+65\%$), remote homology --- the harder, low-identity
regime --- flips from below chance to positive ($-0.023\!\to\!+0.088$), and BixBench saturates
($0.850\!\to\!1.000$). The general axis is flat and the biology axis is monotone increasing over
the same models: the two do not trade off. Figure~\ref{fig:p2traj} plots the joint trajectory,
the two-dimensional form of the capability-shaping curve.

\begin{table*}[t]\centering\small
\caption{Controlled mixture sweep from it-bio ($100$M tokens each, single fixed recipe).
Raising biological share preserves the general axis (MMLU, GSM8K) while monotonically lifting
the biology axis. All values loglikelihood-scored; homology/BixBench are MCC.}
\label{tab:p2mix}
\begin{tabular}{lccccccc}
\toprule
Model & Bio\% & MMLU & GSM8K & TruthQA & Homology-std & Homology-rem & BixBench \\
\midrule
\texttt{M0-base}  & 0  & \textbf{0.802} & \textbf{0.873} & 0.545 & 0.228 & $-0.023$ & 0.850 \\
\texttt{M1-bio5}  & 5  & 0.801 & 0.867 & 0.524 & 0.270 & $-0.055$ & 0.882 \\
\texttt{M2-bio20} & 20 & 0.800 & 0.867 & 0.521 & 0.217 & $-0.074$ & \textbf{1.000} \\
\texttt{M3-bio50} & 50 & 0.796 & 0.861 & 0.519 & \textbf{0.375} & \textbf{0.088} & \textbf{1.000} \\
\bottomrule
\end{tabular}
\end{table*}

\begin{figure}[h]\centering
\includegraphics[width=0.78\linewidth]{figs/p2_fig1_mixture.pdf}
\caption{Capability shaping by biological mixture ratio. As the biological share of a fixed
$100$M-token CPT corpus rises from $0$ to $50\%$, the general axis (MMLU, GSM8K; left scale)
stays flat while the biology axis (homology MCC, BixBench MCC; right scale) rises --- controlled
evidence that mixture ratio shapes the capability profile without a forgetting penalty.}
\label{fig:p2mix}
\end{figure}

\paragraph{Result 3: different scientific data shape different capabilities.}
Holding the biological share at $20\%$ and varying its content (Table~\ref{tab:p2comp},
Figure~\ref{fig:p2comp}) reveals that the \emph{type} of scientific data matters and is not
interchangeable. DNA is most effective for \emph{remote} homology (MCC $+0.043$, versus
$-0.004$ for protein-only), whereas protein sequence is strongest on BixBench ($0.966$) and
standard homology; the protein+DNA mix is intermediate. The composition of a scientific-data
corpus is thus a second, finer lever: not only \emph{how much} scientific data, but \emph{which}
kind, selects which downstream capability is amplified.

\begin{table}[h]\centering\small
\caption{Composition ablation at fixed $20\%$ biological share. Different scientific-data types
amplify different biological capabilities. MCC; \texttt{M2-bio20} = protein+DNA+literature.}
\label{tab:p2comp}
\begin{tabular}{lccc}
\toprule
Model & Homology-std & Homology-rem & BixBench \\
\midrule
\texttt{C1-protein}  & 0.247 & $-0.004$ & \textbf{0.966} \\
\texttt{C2-dna}      & \textbf{0.312} & \textbf{0.043} & 0.850 \\
\texttt{C3-protdna}  & 0.237 & 0.025 & 0.914 \\
\texttt{M2-bio20}\,(+lit) & 0.217 & $-0.074$ & \textbf{1.000} \\
\bottomrule
\end{tabular}
\end{table}

\begin{figure}[h]\centering
\includegraphics[width=0.72\linewidth]{figs/p2_fig2_composition.pdf}
\caption{Composition ablation (biological share fixed at $20\%$). DNA most improves remote
homology; protein most improves BixBench and standard homology --- the type of scientific data,
not just its quantity, selects which capability is shaped.}
\label{fig:p2comp}
\end{figure}

\begin{figure}[h]\centering
\includegraphics[width=0.56\linewidth]{figs/p2_fig3_shaping2d.pdf}
\caption{Two-dimensional capability-shaping trajectory of the mixture sweep. Each point is one
CPT model; increasing biological share moves the model up the biology axis with negligible
displacement along the general axis.}
\label{fig:p2traj}
\end{figure}


% ===================================================================
\section{Part III --- A small-model, matched-compute causal probe}
\label{sec:phase3}

Part II establishes that, at 26B parameters, a biologically-mixed CPT corpus leaves the general
axis essentially flat while lifting biology --- the two capabilities do not trade off, but general
ability is \emph{preserved}, not \emph{improved}. This leaves open a narrower and more provocative
question, motivated directly by the forward-transfer literature this paper sits alongside
\citep{wang2026biopaws}: biological sequence and natural-language paraphrase share a
``difference-detection'' structure --- both require judging whether two strings that look
superficially similar (or dissimilar) are functionally equivalent. If that shared structure is
real, biological CPT should \emph{measurably improve} the natural-language skills that are its
direct structural counterpart, not merely leave them undisturbed. Testing this requires a regime
with capability headroom, which the near-ceiling 26B model does not offer on these metrics; we
therefore turn to a small model, where headroom is available and any change is easier to attribute.

\paragraph{A causal, matched-compute design.} We start from GPT-2 (124M parameters,
\citealp{radford2019gpt2}) and continue pretraining it on English web text (openwebtext) for
$300$M tokens, producing a single shared checkpoint (\texttt{nl\_base}). From this \emph{one}
checkpoint we then run two arms for exactly $200$M further tokens each ($3{,}051$ optimizer
steps, identical batch size, learning rate, and schedule) --- the only difference is what the
extra tokens contain: the \textbf{NL-arm} continues on openwebtext; the \textbf{Bio-arm}
continues on a $50$/$50$ protein/DNA mixture. Both arms are full-parameter finetunes (no LoRA;
at 124M this is cheap and avoids any confound from adapter capacity). We repeat the two-arm
comparison at three random seeds (data order and optimizer), giving a mean and standard deviation
for every metric rather than a single run. Because both arms share the identical start, token
budget, and optimizer trajectory, any difference between them isolates \emph{what the extra
pretraining tokens contained}, not how much training occurred.

\paragraph{Two families of NL probes.} We evaluate both arms zero-shot on: (i) a \emph{structural/
association} family --- PAWS-en \citep{yang2019pawsx}, the direct English counterpart of the
paraphrase-detection task whose zero-shot transfer to protein homology motivated this line of
work, plus the five-task character-manipulation battery from \citet{brown2020gpt3} (anagram
solving on the middle letters, letter-cycling, random-letter insertion, and word-reversal), which
probes sub-word structure the way homology probes sub-sequence structure; and (ii) a
\emph{knowledge/discourse} family --- HellaSwag and LAMBADA --- included as a control where we
expect biological data to be neutral or mildly harmful, since these tasks reward broad
world-knowledge and long natural-language discourse rather than structural matching.

\begin{table}[h]\centering\small
\caption{Small-model matched-compute continuation (GPT-2, $200$M tokens/arm, mean over 3 seeds,
identical start/budget/optimizer in both arms). PAWS-en shows a reproducible biological lift;
the character battery is floored for both arms at this scale (not informative); the
knowledge/discourse family shows the expected cost, sharpest on long-range discourse (LAMBADA).}
\label{tab:p3}
\begin{tabular}{llccc}
\toprule
Family & Task & NL-arm & Bio-arm & $\Delta$ (Bio$-$NL), mean$\pm$sd \\
\midrule
Structural & PAWS-en                  & 0.498 & \textbf{0.529} & $\mathbf{+0.031\pm0.009}$ \\
Structural & Anagrams (mid, 1/2)      & 0.000 & 0.000          & $0.000$ (floor, both arms) \\
Structural & Cycle letters            & 0.000 & 0.000          & $0.000$ (floor, both arms) \\
Structural & Random insertion         & 0.000 & 0.000          & $0.000$ (floor, both arms) \\
Structural & Reversed words           & 0.000 & 0.000          & $0.000$ (floor, both arms) \\
Knowledge  & HellaSwag                & 0.291 & 0.285          & $-0.006\pm0.000$ \\
Knowledge  & LAMBADA (discourse)      & 0.333 & 0.188          & $-0.146\pm0.004$ \\
\bottomrule
\end{tabular}
\end{table}

\paragraph{Result: a real, reproducible, but narrow transfer.} Table~\ref{tab:p3} shows a single
clean positive result rather than a uniform one. PAWS-en accuracy rises by $+3.1$ points
($0.498\!\to\!0.529$) in the Bio-arm, and the sign is identical across all three seeds
(sd $0.009$, an order of magnitude smaller than the effect) --- to our knowledge the first
demonstration that continuing pretraining on \emph{raw biological sequence}, with no natural-
language paraphrase data in the mix, measurably improves a model's ability to detect natural-
language paraphrase, under a design that holds compute exactly fixed between arms. The five
character-manipulation tasks are uninformative at this scale: both arms score exactly $0$,
consistent with \citet{brown2020gpt3}'s own finding that these tasks are near the noise floor
below the tens-of-billions-of-parameters range; a $124$M model, whether or not it has seen
biological sequence, simply cannot do them, so this comparison awaits a larger base model. On the
knowledge/discourse control, the cost is real: HellaSwag drops slightly ($-0.6$ pt) and LAMBADA
drops sharply ($-14.6$ pt). We read the LAMBADA drop as the visible signature of a genuine
trade-off at this scale that Part II's 26B MoE model does not show: a $124$M model trained fully
(not via LoRA) and pushed with a $50\%$-biological continuation has limited spare capacity, so a
half-biological corpus measurably crowds out long-range natural-language discourse modeling even
as it sharpens a narrower structural skill. The dissociation is therefore partial, not clean: one
structural task transfers positively and reproducibly, one discourse task pays a real cost, and
the remaining structural probes are simply unmeasurable at this scale --- an honest result that
neither confirms nor refutes a fully general ``biology helps all structure'' claim, but does
establish, for the first time under a causal matched-compute design, that the transfer is real
for at least one task in the family motivated by this paper's overall thesis.


% ===================================================================
\section{Discussion}
\label{sec:discussion}

\subsection{A division of labor between CPT and SFT}
The recurring shape of our results --- base $<$ BioCPT $>$ SFT on every held-out axis --- points
to a division of labor that is usually conflated. Continued pretraining does not simply add
domain facts; it \emph{re-organizes and lifts a shared capability substrate} that is then
visible on tasks (general knowledge, code) far from the pretraining domain. Supervised
fine-tuning does the opposite: it \emph{narrows} the model onto the target task distribution,
trading breadth on held-out axes for focus on the served task. Neither stage is a substitute
for the other, and budgeting them as interchangeable --- or omitting CPT to ``save compute'' ---
misreads their roles. This is consistent with, and sharpens, prior router-level evidence on this
model family that CPT reshapes middle-layer computation while SFT re-aligns the output.

\subsection{Biological sequence as structured scientific data}
If narrow, out-of-distribution biological sequence \emph{raises} general and coding competence,
then the ``narrow domain = capacity competition'' framing is inadequate. We propose instead that
biological sequence functions as \emph{structured scientific data}: a corpus whose value to a
foundation model lies not (only) in the facts it carries but in the structural regularities it
forces the model to internalize --- regularities that transfer to other structured/formal
languages such as code. Under this view biology is simply the first representative member of a
broader class (code, mathematics, chemistry/SMILES), and the natural next question is not
``does biology help biology'' but ``how does the \emph{composition} of scientific-data CPT
shape the capability profile of a foundation model.''

\subsection{Mixture ratio is a controllable, general-preserving lever}
The controlled sweep of \S\ref{sec:phase2} answers that question for this model. The two axes do
not trade off: across a biological share of $0\!-\!50\%$ the general axis is flat while the
biology axis rises monotonically, so there is no ``forgetting threshold'' to budget around
within this range. Equally, the composition ablation shows the lever has structure --- DNA and
protein amplify \emph{different} biological capabilities --- so a practitioner can, in principle,
target a desired capability profile by choosing not just how much scientific data to add but
which kind. This is the operational form of the data-centric thesis: training-data composition
is a design surface for the capability profile, not merely a source of facts.

\subsection{Limitations}
This is a training-free re-analysis of one lineage of one 26B MoE model; the CPT / SFT recipe is
fixed, so we characterize \emph{this} lineage rather than establishing a scaling law. Part of the
ARC/HellaSwag gain is a format effect (instruct-vs-completion loglikelihood), which we flag and
which does not affect the knowledge-heavy MMLU gain. Zero-shot sequence-homology is not
measurable by simple probes, so that axis is argued indirectly. TruthfulQA is a genuine, if
small, regression and we do not hide it. The controlled experiment of \S\ref{sec:phase2}
addresses the observational limitation directly --- it fixes the recipe and varies only the data
mixture --- but is itself bounded: each model sees $100$M CPT tokens (a pilot budget chosen to
find trends, not a scaling law), on one 26B MoE model. Whether the general-preserving,
biology-lifting mixture behavior holds at $10\times$ the token budget, across model sizes, and
for other scientific-data domains (code, mathematics, chemistry) remains open, and is the natural
next step now that the lever itself is established.


\section*{Data and code availability}
All three (four) checkpoints derive from a public lineage of Gemma-4-26B-A4B; the merged
BioCPT checkpoint is released as \texttt{dnagpt/OmniGene-4-CPT-v2-merged}. Evaluation scripts
(general, coding, biology-loglikelihood, and chain-of-thought diagnostics), the exact few-shot
configurations, and per-example model outputs for every table in this paper are released at
the project repository. All benchmarks (MMLU, ARC, HellaSwag, TruthfulQA, MBPP, HumanEval,
GSM8K, BixBench) are public; no proprietary data were used. The Part~I re-analysis is
training-free; the seven controlled Part~II CPT models ($100$M tokens each) together with their
QLoRA adapters, data-mixture recipes, and per-example evaluation outputs are also released at
the project repository. Part~II reproduces on a single 96\,GB GPU. The Part~III GPT-2 models
(one shared \texttt{nl\_base} plus six matched-compute continuation checkpoints, $3$ seeds
$\times$ 2 arms) are full-parameter finetunes released at the same repository together with the
character-manipulation probe data (mirrored from the original GPT-3-paper release) and reproduce
on a single consumer-grade GPU in under two hours total.

\section*{Acknowledgments}
% [placeholder]

\section*{Declaration of interests}
The author declares no competing interests.

\bibliographystyle{unsrtnat}
\bibliography{refs}

\end{document}
